\documentclass[10pt]{article}
\usepackage[margin=20mm]{geometry}
\usepackage{amsmath,amssymb,fontspec,unicode-math,xurl,hyperref}
\IfFontExistsTF{DejaVu Serif}{\setmainfont{DejaVu Serif}}{\setmainfont{Latin Modern Roman}}
\setmathfont{latinmodern-math.otf}
\setlength{\parindent}{0pt}\setlength{\parskip}{6pt}\sloppy
\begin{document}
ELVIA Scientific Verification Report
0100 / I25-43469\par
Towards LLM Unlearning Resilient to Relearning Attacks: A Sharpness-Aware Minimization Perspective and Beyond\par
Inspection outcome: Major. The implication from SAM stationarity to vanishing directional curvature is false under the stated smoothness and nonzero-gradient conditions. Complete counterexamples demonstrate a substantial invalid derivation tied to a principal contribution. The reported relative empirical improvements are not refuted.\par
Additional outcomes: Minor for a bounded rho-ablation interpretation and for concrete public-release defects; Needs Evidence for linkage to published runs and inconsistent settings/values. Scientific decisions and controller acceptance remain pending. Author adoption has not been requested.\par
Batch B001 / [internal task]; token [internal identifier omitted]. Review date: 24 October 2026. AI-assisted verification; no human-review approval is asserted.\par
\# 1. Identity, versions and scope\par
Authors: Chongyu Fan, Jinghan Jia, Yihua Zhang, Anil Ramakrishna, Mingyi Hong and Sijia Liu. Exact source: PMLR 267 (ICML 2025), pp. 15762-15778, publisher-linked fan25e.pdf, 17 pages [R1]. Locations use one-based PDF pages, equal to printed article pages 1-17, rather than volume pagination.\par
The existing papers/0100/0100.pdf was identified and SHA-256 matched to verification/0100/controller\_acquisition\_receipt.json before analysis; its bytes were preserved. SHA-256: 77fc07b5fbaff3f32605be16547ac8e131441ada2a05c0e7bb04f8f866a34c29.\par
The latest arXiv version observed was 2502.05374v4, 27 May 2025 [R2], a preconference comparison rather than a postpublication correction. Its p. 4 retains the affected inference. The official repository was pinned at ea25a74b509b25ece06e65d59bbff01e6d2239da (27 September 2025); 40 selected source/configuration files were acquired preserving their directory tree and verified against Git blobs [R3]. Relevant expressions are unchanged from the first-code tree, a4ec8e730052b2014c8a5d4875b11ec18dcd374c (23 February 2025). Release history does not identify publication-generating code.\par
Checked: Eqs. (1)-(12), Algorithm A1, main results, appendices A-F, metric directions, table arithmetic, threat descriptions, selected training/evaluation paths and targeted version comparisons. Calculations are analytic, symbolic or finite reviewer checks. No author code was run/imported; no model, training, attack, paid API or empirical reproduction was used. Raw trials, checkpoints and publication-run mappings were not obtained. Coverage and exclusions are itemized in coverage.json and the specialist audits [R6].\par
\# 2. Definitions and valid reasoning retained\par
Write ell for the forget loss, g for its gradient, H for its Hessian, and v=g/||g|| where g is nonzero, as explicitly assumed on p. 4. The robust formulation restricts weight perturbations to a norm ball, with default Euclidean radius rho>0, and relates relearning to an objective counteracting the forget objective. This is a conditional threat model, not an equivalence to every fine-tuning attack.\par
Equation (4) explicitly linearizes the inner maximization; rho*v correctly optimizes that linearized problem. Equation (6) is explicitly a first-order Taylor approximation. Equation (7), gradient of the gradient norm equals H*v, is correct for nonzero g. Equation (8) is a local finite-difference approximation. Equations (9)-(11) define coherent smoothing/regularization objectives under ordinary integrability/differentiability conditions. These valid steps survive.\par
Algorithm A1 perturbs the forget loss, computes its gradient at perturbed weights, and combines it with the retain gradient at original weights. The public SAM paths follow this ordering: WMDP base.py lines 764-775, restoration/retain backward at 834-849, combination at 853-864; MUSE sam\_iterative.py lines 129-138 and helpers through 226. This positive static check assumes nonzero forget gradients, relevant gradient entries for trainable parameters and one accumulation step, explicitly configured in WMDP. General multi-microbatch accumulation, AMP and distributed equivalence were not certified [R3,R6].\par
WMDP UE is 1 minus benchmark accuracy, so higher is better; UT is MMLU accuracy. MUSE forget-set VerbMem and KnowMem are lower-is-better; retain-set KnowMem is higher-is-better. Relative resistance, utility, total information erasure and universal attack robustness are distinct claims.\par
\# 3. Problem 1 - Stationarity does not force curvature to vanish\par
Source claim and location. Section 4, p. 4, immediately after Eq. (8), concludes from solving Eq. (5) toward stationarity that ||H*v|| tends to zero. This is the mathematical bridge from SAM optimization to curvature suppression and the subsequent CR/smoothing motivation, closely tied to a principal contribution [R1,R2].\par
Suspect step. A zero gradient of the full surrogate does not force one component to vanish. For the Eq. (6) surrogate G, stationarity permits cancellation between g and rho*H*v. For F in Eq. (5), the derivative contains the shifted gradient and the Jacobian of v. The correct identities appear below. Adding a retain gradient cannot justify dropping the other terms.\par
Complete counterexample. Fix any rho>0, for example 0.01, and take the scalar polynomial ell displayed below. Choose lambda=0, permitted by the nonnegative retain coefficient, or a constant retain loss with any lambda. The loss is smooth, coercive and globally nonnegative. Its only critical points are -rho, rho/2 and rho, with values 5*rho\textasciicircum{}2/12, 215*rho\textasciicircum{}2/192 and 13*rho\textasciicircum{}2/12, so its global minimum is positive.\par
The factored gradient is positive on (-rho,rho/2), hence v=1 near x=0. The source’s nonzero-gradient condition holds at and near the solution. At zero, g=rho/2, H=-1 and ||H*v||=1. Locally Eq. (5) is exactly F(x)=ell(x+rho). Thus F’(0)=ell’(rho)=0 and F”(0)=ell”(rho)=1>0: zero is a strict local minimum of the stated surrogate, but directional curvature does not vanish. No undefined normalization or stationary maximum is used.\par
Independent verification. With rho=0.01, step size 0.25 and x0=0.001, 80 reviewer scalar iterations yield x=7.45278e-14, g=0.00499999999993 and ||H*v||=1.00000000000745. At zero, the Eq. (8) finite difference with mu=1e-6 is -1.00004999, tending to -1. Exact symbolic checks and the complete replay script are preserved [R6].\par
Alternative readings and counterevidence. This tests Eq. (5) as the paper’s approximate-SAM surrogate; it does not claim its perturbation is an exact nonlinear inner maximizer. Algorithm A1’s deliberate omission of higher derivatives cannot remove the example because v is locally constant. Appendix A supplies a second smooth nonnegative loss with strict local minima for both Eq. (5) and Eq. (6), while ||H*v||=1, closing the alternative interpretation of the paragraph. These are generic mathematical witnesses, not simulated LLM losses or NPO training runs.\par
Impact and repair. Major: a core-connected inference is false even if the experiments and broader smoothing hypothesis survive. Valid local approximations and conditional robust optimization remain available. The derivative identities below are verified scoped replacements; no complete argument restoring the claimed zero-curvature conclusion has been verified.\par
\[g=\nabla\ell,\quad v=g/\|g\|_2,\quad Dv=(I-vv^\top)H/\|g\|_2.\] \[\begin{aligned}F(\theta)&=\ell(\theta+\rho v(\theta)),&
\nabla F&=(I+\rho Dv)^\top\nabla\ell(\theta+\rho v),\\
G(\theta)&=\ell(\theta)+\rho\|g\|_2,&
\nabla G&=g+\rho Hv.\end{aligned}\] \[\begin{aligned}\ell(x)&=\rho^2+\frac{x^4}{4\rho^2}-\frac{x^3}{6\rho}-\frac{x^2}{2}+\frac{\rho x}{2},\\
\ell'(x)&=\frac{(x-\rho/\allowbreak{}2)(x-\rho)(x+\rho)}{\rho^2}.\end{aligned}\]\par
\# 4. Problem 2 - The large-radius ablation confuses forgetting with utility\par
Source claim. The rho-ablation discussion on p. 9 and Appendix E, p. 15, attributes rho=0.1 to ineffective forgetting. Table A3 instead shows identical pre-attack forget scores, substantial utility deterioration and mixed post-attack changes [R1,R6].\par
Complete check. Pre-attack forget VerbMem and KnowMem are both 0.00 at both radii. Increasing rho from 0.01 to 0.1 reduces retain KnowMem by 11.57 points; after attack, forget VerbMem improves by 5.42 while KnowMem worsens by 2.19. Directions follow the paper’s explicit UE/UT definitions.\par
Alternative reading and impact. If effectiveness informally means the joint forget/retain objective, the utility loss supports a worse tradeoff; the specific assertion about preventing forgetting does not follow. Minor: a bounded interpretation error, with no demonstrated failure of the selected rho=0.01 comparison. A verified replacement is: pre-attack forgetting is unchanged at displayed precision, retain utility falls substantially, and post-attack metrics show a mixed tradeoff. This repairs the description, without asserting new experimental results.\par
Table A3 metric | rho=0.01 | rho=0.1 | Change\par
Retain KnowMem (higher) | 35.48 | 23.91 | -11.57; worse\par
Pre-attack VerbMem / KnowMem (lower) | 0.00 / 0.00 | 0.00 / 0.00 | No displayed change\par
Post-attack VerbMem (lower) | 58.38 | 52.96 | -5.42; better\par
Post-attack KnowMem (lower) | 38.33 | 40.52 | +2.19; worse\par
\# 5. Problem 3 - The public GP penalty is detached\par
Source claim and location. Equation (10), p. 5, defines a differentiable gradient-norm penalty. The public WMDP path is NPO\_FT\_GNR, WMDP/src/unlearn/GA.py lines 218-264, reachable through \_\_init\_\_.py lines 24-25 and the supplied GP command [R3].\par
Complete static check. Lines 255-257 request differentiable gradients with create\_graph=True, but line 258 converts squared-gradient sums using .item(). The accumulated norm and its square root are Python scalar values. Adding that value at line 262 contributes zero parameter derivative. The official Tensor.item documentation confirms nondifferentiability [R5]. For the reviewer example ell(theta)=theta\textasciicircum{}2/2 at theta=1 and rho=0.01, the intended derivative of ell+rho*|ell’| is 1.01; a detached penalty leaves 1.0. No author or Torch training code was executed.\par
The command supplies gnr\_rho=10, while this loss hardcodes 0.01. A changed logged loss is distinct from a changed optimization gradient. The expression is present in the first code release, so it is not only a recent edit.\par
Impact and counterevidence. Minor for the demonstrated release defect; published GP-result implications remain Needs Evidence because run-source linkage is missing. Objective identity would matter even if scores remained favorable after such linkage were established. Preserving a differentiable tensor norm and using the intended coefficient are necessary local requirements, but no complete repaired trainer or empirical restoration was verified.\par
\# 6. Problem 4 - The supplied CR route has no dispatch branch\par
Source claim and location. Equation (11), p. 5, defines curvature regularization. WMDP/run\_wmdp\_unlearn.sh lines 38-48 selects NPO+FT+CR, but the complete dispatch function in WMDP/src/unlearn/\_\_init\_\_.py lines 5-28 has no branch for that name and falls through to ValueError. Importing its class does not make the command reach it [R3,R6].\par
Separately, the stored NPO\_FT\_CR class, GA.py lines 267-346, calls backward inside compute\_loss at line 299 and detaches the gradient-difference norm with .item() at line 340, adding a scalar at 344. Adding dispatch alone would not validate the penalty or gradient accounting.\par
Impact and counterevidence. Minor for this public release path; published-result impact Needs Evidence. The unreachable command is counterevidence against asserting that this class certainly generated the reported CR measurements. Another execution path remains possible. No command was run, no training failure was reproduced, and no complete repair is claimed.\par
\# 7. Problem 5 - The released MUSE objective is not likelihood-ratio NPO\par
Source claim and location. The paper identifies NPO as the MUSE baseline on pp. 3 and 7-8. The original NPO definition depends on the conditional probability of a forget response relative to a reference, Eq. (3) in Zhang et al. [R4]. The saved MUSE/baselines/baselines/iterative.py lines 155-157 and sam\_iterative.py lines 256-258 instead apply logsigmoid to reference.logits-current.logits over the full tensor, without probability normalization or selecting the forget labels for that term [R3,R6].\par
Complete invariant check. Set beta=0.1, reference and initial current logits (1,-1), and subsequent current logits (-1,-3), a common shift of -2. All logits are finite, beta is positive, and the reference and forget label remain fixed. Both current softmax distributions equal (0.88079707798,0.11920292202), hence every response likelihood is unchanged. Proper likelihood-ratio NPO remains 13.8629436112; the released loss changes from 13.8629436112 to 11.9627773876. This establishes objective mismatch without assuming a real model follows a common-offset trajectory.\par
Counterevidence and impact. Both MUSE baseline and SAM use the expression, so a relative advantage can survive. WMDP uses cross-entropy loss differences and is outside this raw-logit counterexample; its reduction convention is a separate scope question. Minor for the released objective mismatch; actual published MUSE objective remains Needs Evidence. If the released expression generated the named NPO comparison, that identity failure would need assessment without requiring score deterioration. No run mapping or complete replacement with verified label shifting, masking, sequence reduction and beta convention was obtained.\par
\[L_{\rm NPO}=-\frac{2}{\beta}\mathbb E\log\sigma\!\left(-\beta\log\frac{\pi_\theta(y\mid x)}{\pi_{\rm ref}(y\mid x)}\right),\qquad \beta>0.\]\par
\# 8. Empirical checks and unresolved evidence\par
Reported arithmetic supports relative robustness in the tested settings. Every smoothing variant in Table 2 exceeds NPO in all displayed attack columns. SAM-minus-NPO UE gains are 0.13/0.11/0.08 at 20/40/60 samples and 0.13/0.23/0.22 at 1/2/3 epochs. Table A1 improves every displayed attacked entry for NPO, GradDiff and RMU with SAM. Table A2 reduces post-attack VerbMem/KnowMem by 5.10/2.84 points for News and 9.14/7.00 for Books. These verify published arithmetic, not independent experiments [R1,R6].\par
Tradeoffs survive alongside the gains. Table 2 UT is 0.42 for NPO+SAM versus 0.44 for NPO; SAM UE falls from 0.74 before attack to 0.45 after the 60-sample attack. Table 1 gives identical SAM/TAR UE (0.74 before, 0.70 after), but SAM UT=0.51 versus TAR=0.54. Reported runtimes yield approximately 647.12 times TAR/SAM and 1.98 times SAM/NPO; hardware-normalized timings were not reproduced. Favorable resistance and efficiency do not imply immunity or dominance on all metrics.\par
Figure 5’s p. 8 caption lists Books 200/300/400 and News 400/500/600 samples; neighboring prose says 75/100/125. Released commands use Books 50/75/100 steps and News 100/125/150 with batch size 4. With one process, no accumulation and full batches, processed-example counts match the caption; distinct-sample counts and publication settings remain unresolved. Figure 2a gives NPO+WA Relearn60 UE=0.42, while Table 2 gives 0.40. Both exceed NPO=0.37 and remain below SAM=0.45, preserving that ordering. Original run settings and aggregates are needed.\par
The Books SAM shell command uses retain coefficient 2.5, outside Appendix C’s stated [0.25,1.0] search range. WMDP shell overrides mean JSON defaults alone cannot identify final settings. These are provenance/configuration concerns, not demonstrated changes in published scores. Trial outputs, training-seed uncertainty and complete selected configurations remain unavailable.\par
The README’s enhanced-GCG command is parameterized by the chosen unlearned model and specifies optimization. This supports intended per-model attack adaptation and is counterevidence to a transferred-prefix-only accusation. Actual prefixes, seeds, external attack revision and logs were not verified. Two-dimensional loss landscapes and token-distribution changes are evidence in the displayed settings, not certificates of all-direction curvature suppression or total information erasure. No new attack requirement has been imposed.\par
Public issue \#1 asked about MUSE GPU count and model availability. Collaborator a-F1 replied on 17 June 2025 that experiments used eight A6000 GPUs; it closed on 18 June. The all-states issue endpoint returned only this issue, and both comments were inspected [R7]. This is historical hardware clarification, not a correction of Problems 1 or 3-5. It does not specify data-parallel world size versus model sharding (the release uses device\_map=’auto’), accumulation or run mapping; multiplying step-times-batch counts by eight is therefore not justified. Private/external correspondence was not assessed.\par
Other observations were not promoted to substantive findings: Eq. (12)’s averaging index and p. 6’s printed prediction-loss sign appear inconsistent with surrounding descriptions, but the evidence does not establish resulting numerical damage or explain the plots.\par
\# 9. Dependence, version status and conclusions\par
The Major judgment rests on Problem 1 alone: a substantial invalid inference in the central SAM-to-curvature explanation. It does not depend on a quota, release discrepancies, difficulty of repair or empirical-score failure. Curvature involvement in differentiation is established; convergence to zero directional curvature is not. The theoretical claim and empirical observations have separate evidentiary status.\par
The observed preconference arXiv v4 retains the inference. Repository history preserves the release defects from its first code release; the public hardware clarification does not repair them or establish which source generated published scores. No complete theoretical/implementation repair or author adoption is verified. Determining actual GP/CR/NPO experimental implications requires original run-source/checkpoint records and, if corrected objectives change the evaluated method, fresh experiments.\par
Concern / outcome | Supported impact | Repair verification | Adoption / remaining work\par
Problem 1: analytic counterexamples | Major; central derivation invalid | Correct derivative identities; guarantee not restored | Not requested; valid assumptions/proof needed\par
Problem 2: rho interpretation | Minor; local interpretation | Replacement interpretation verified by table | Not requested; distinguish UE and UT\par
Problems 3-4: GP / CR release | Minor release defects; published impact unresolved | No complete trainer repair | Run linkage; objective/gradient validation\par
Problem 5: MUSE NPO release | Minor mismatch; published impact unresolved | Invariant check verified; no full replacement | Run linkage; possibly fresh experiments\par
Budget / WA / configuration | Needs Evidence | No corrected value or run identified | Original configs, trial aggregates and checkpoints\par
SAM order / relative scores | No Issue Found within inspected scope | Static ordering and arithmetic only | No empirical reproduction; decision pending\par
\# Appendix A. A witness also covering Eq. (6)\par
Let t=x/rho, h(t)=1+3t+(9/2)t\textasciicircum{}2+7t\textasciicircum{}3+100t\textasciicircum{}4, q(t)=(t-1)\textasciicircum{}2*h(t), and ell(x)=rho\textasciicircum{}2*q(x/rho). The identity h(t)=(1+3t/2)\textasciicircum{}2+t\textasciicircum{}2*(9/4+7t+100t\textasciicircum{}2) proves h>0: the quadratic has positive leading coefficient and discriminant -851. Thus ell is smooth and nonnegative.\par
Exact derivatives: q’(0)=1, q”(0)=-1, q”’(0)=6, q’(1)=0, q”(1)=231. At x=0, g=rho>0, H=-1 and v=1 nearby. The Eq. (5) surrogate ell(x+rho) has derivative zero and second derivative 231. The Eq. (6) surrogate ell(x)+rho*ell’(x) has derivative rho*(1-1)=0 and second derivative -1+6=5. Both have strict local minima while ||H*v||=1. Constant retain loss leaves the calculation unchanged. This closes the Eq. (6) reading without invoking a zero base gradient [R6].\par
\# Appendix B. Replay, hashes, failures and roles\par
Every supporting file is inside sources/0100/ or verification/0100/. Main ledgers: math\_audit.json, empirical\_audit.json, static\_code\_audit.json and public\_issue\_audit.json. Reviewer scripts/results include curvature\_check\_verified.py and curvature\_check\_results.json, empirical\_independent\_checks.py/json, static\_source\_checks.py and static\_code\_checks.json. Extracted text, numbered code, repository trees/history, acquisition manifests and source renders are retained. artifact\_manifest.json records required-file hashes; source identity remains linked to the controller receipt.\par
Actual reviewer computation used existing Python/SymPy for algebra, standard arithmetic for score differences and loss invariance, and static text/blob checks. The initial curvature\_check.py failed on SymPy structural equality of equivalent factorizations. It and curvature\_check\_initial\_failure.json were preserved; the verified script simplifies expression differences and completed successfully. The first Tensor.item download was an HTML redirect, preserved before following its explicit versioned target. Neither HTML document was executed. A report-construction call also had a JavaScript quoting syntax error before execution; no file was modified by that failed call, and the construction was split into valid calls.\par
Additional SHA-256 values: arXiv v4 PDF e2bea9fb34f3e182a47c5be69d6b4cfd25bb73f4848c611f64f221885792aba5; NPO-definition HTML 0058376a7d11cf2adcf325f25c2a51d03fe7b1a40cf9dfa75a653da2f7ae6b8e; Tensor.item 2.14 HTML 622422398940e918a848581687a56f1f2a885b8df1665799e780f08a1a8577c1. NPO is consulted only for its objective definition, not as another reviewed paper.\par
Root AI verified identity, integrated evidence, inspected source/report renderings and assembled deliverables. Child AI audit\_claims\_0096 checked 0100 mathematics; audit\_code\_0096 inspected its code/history; capacity\_0098 checked its empirical values and figures. No human approval, author adoption, author message, publication, software installation, paid experiment or queue/controller change is asserted or performed.\par
The proposed report PDF is generated by compiling the self-contained TeX with the existing offline compiler. The acceptance record identifies the actual compilation and PDF checks for this version; scientific decisions and controller promotion remain pending. Equivalent reviewer-only replay commands, from the assigned staging root, are:\par
{\ttfamily\footnotesize \& '[local archive path omitted]' `\\
 'verification\textbackslash{}0100\textbackslash{}curvature\_check\_verified.py'\\
\& '[local archive path omitted]' `\\
 'verification\textbackslash{}0100\textbackslash{}empirical\_independent\_checks.py'\\
\& '[local archive path omitted]' `\\
 'verification\textbackslash{}0100\textbackslash{}static\_source\_checks.py'}\par
\# Consulted references and evidence identifiers\par
1. Fan et al. ICML 2025 publisher version, PMLR 267:15762-15778. Local source papers/0100/0100.pdf; receipt and identity check in verification/0100/. https://proceedings.mlr.press/v267/fan25e.html\par
2. Fan et al. arXiv:2502.05374v4, 27 May 2025. PDF/HTML/extraction in sources/0100/; version\_acquisition.json records retrieval; targeted p. 4 comparison. https://arxiv.org/abs/2502.05374v4\par
3. OPTML-Group/Unlearn-Smooth, commit ea25a74b509b25ece06e65d59bbff01e6d2239da; first code a4ec8e730052b2014c8a5d4875b11ec18dcd374c. Saved tree/files/history and manifests in sources/0100/. Line references use the pinned release. https://github.com/OPTML-Group/Unlearn-Smooth/tree/ea25a74b509b25ece06e65d59bbff01e6d2239da\par
4. Zhang, Lin, Bai and Mei. Negative Preference Optimization: From Catastrophic Collapse to Effective Unlearning. arXiv:2404.05868v2, 10 October 2024, Section 3 Eq. (3). Definition-only support, saved npo\_definition\_2404\_05868v2.html. https://arxiv.org/html/2404.05868v2\par
5. PyTorch Tensor.item, version 2.14 API documentation, retrieved 24 October 2026. Supports scalar conversion/nondifferentiability; not asserted to be the publication training environment. https://docs.pytorch.org/docs/2.14/generated/torch.Tensor.item.html\par
6. Reviewer evidence, 24 October 2026: verification/0100/math\_audit.json, empirical\_audit.json, static\_code\_audit.json, exact scripts/results, source renders and provenance.json. Analytic, static and finite checks; no author execution.\par
7. Official repository public issue \#1 and both comments. Opened 9 April 2025; collaborator response 17 June; closed 18 June. Hardware/model-availability clarification only. API responses and public\_issue\_audit.json preserved. https://github.com/OPTML-Group/Unlearn-Smooth/issues/1\par
\end{document}
